\section*{Hoofdstuk \ref{ch:b1:e-ph-diagrams} --- E-pH-diagrammen}

\begin{solution}{exo:b1:e-ph-diagrams:1}
Mangaan: \ce{Mn} (0) $<$ \ce{Mn^2+}, \ce{Mn(OH)2} ($+$II) $<$ \ce{MnO2}
($+$IV) $<$ \ce{MnO4-} ($+$VII). Chloor: \ce{Cl-} ($-$I) $<$ \ce{Cl2} (0)
$<$ \ce{HClO}, \ce{ClO-} ($+$I) $<$ \ce{ClO3-} ($+$V).
\end{solution}

\begin{solution}{exo:b1:e-ph-diagrams:2}
$-0.059 \times 3/1 = -0.177$; $+0.059 \times 2/2 = +0.059$;
$-0.059 \times 8/5 = -0.094$; $-0.059 \times 2/2 = -0.059$ V per
pH-eenheid. Voor het koperkoppel komen de protonen vrij aan de kant van
de gereduceerde vorm ($m = -2$): de potentiaal stijgt met de pH.
\end{solution}

\begin{solution}{exo:b1:e-ph-diagrams:3}
\ce{H+}/\ce{H2} gaat bij pH 0 door 0~V en bereikt 0.50~V nooit voor
$\mathrm{pH} \geqslant 0$. \ce{O2}/\ce{H2O} bereikt 0.50~V bij
$\mathrm{pH} = 0.73/0.059 = 12.4$ en 0~V pas bij pH 20.8, buiten de
schaal. De band is bij elke pH \qty{1.23}{V} breed, ook bij pH 7 (van
$-0.41$ tot $+0.82$~V).
\end{solution}

\begin{solution}{exo:b1:e-ph-diagrams:4}
(1, 1.0~V): \ce{Fe^3+}. (4, 0): de lijn \ce{Fe(OH)3}/\ce{Fe^2+} ligt op
$1.09 - 0.71 = 0.38$~V en het punt ligt eronder, boven $-0.47$~V:
\ce{Fe^2+}. (10, $-0.4$~V): tussen $-0.07 - 0.59 = -0.66$ en $0.28 -
0.59 = -0.31$~V: \ce{Fe(OH)2}. (10, 0.5~V): \ce{Fe(OH)3}.
\end{solution}

\begin{solution}{exo:b1:e-ph-diagrams:5}
\ce{Fe^2+}/\ce{Fe}: $-0.41 + 0.030 \times (-6) = -0.59$~V.
\ce{Fe^3+}/\ce{Fe(OH)3}: $14.00 - \frac13(38.55 - 6) = 3.15$. De gebieden
van de opgeloste ionen (corrosie) worden groter: \ce{Fe^2+} reikt tot
$-0.59$~V en \ce{Fe^3+} tot pH 3.15; metaal en hydroxiden verliezen
terrein.
\end{solution}

\begin{solution}{exo:b1:e-ph-diagrams:6}
\ce{Fe(OH)3 + 3H+ + e- -> Fe^2+ + 3H2O}: $E = E^\circ + 0.059\log\frac{[\ce{H+}]^3}{[\ce{Fe^2+}]}
= (E^\circ - 0.059\log c) - 0.177\,\mathrm{pH}$. Bij pH 1.82 snijdt de
lijn \ce{Fe^3+}/\ce{Fe^2+} op 0.77~V, dus de constante is
$0.77 + 0.177 \times 1.82
= 1.09$~V.
\end{solution}

\begin{solution}{exo:b1:e-ph-diagrams:7}
$E^\circ(\ce{Cu+}/\ce{Cu}) = 0.52 > E^\circ(\ce{Cu^2+}/\ce{Cu+}) = 0.16$:
volgens \cref{prop:b1:e-ph-diagrams:disproportionation-criterion}
disproportioneert \ce{Cu+}. Enige grens: $E = 0.34 + 0.030\log 0.010 =
0.28$~V.
\end{solution}

\begin{solution}{exo:b1:e-ph-diagrams:8}
Bij pH 0 overlapt het gebied van koper (onder 0.28~V) met dat van water
en ligt het boven de waterstoflijn: geen reactie met zoutzuur zonder
lucht. Salpeterzuur: \ce{NO3-}/\ce{NO} op 0.96~V ligt boven het gebied
van koper, dat geoxideerd wordt:
\ce{3Cu + 2NO3- + 8H+ -> 3Cu^2+ + 2NO + 4H2O}.
% ledger: dfg:NO3-_ao, dfg:NO_g (E0 computed in figdata/bachelor-1/standard-potentials.py)
\end{solution}


\begin{solution}{exo:b1:e-ph-diagrams:9}
Bij pH 3 ligt de lijn \ce{O2}/\ce{H2O} op $1.23 - 0.18 = 1.05$~V, boven
het hele gebied van \ce{Fe^2+}: zuurstof oxideert ijzer(II). Boven pH
1.82 slaat het gevormde ijzer(III) neer als het geelbruine hydroxide:
\ce{4Fe^2+ + O2 + 10H2O -> 4Fe(OH)3 + 8H+}.
\end{solution}

\begin{solution}{exo:b1:e-ph-diagrams:10}
Verticale grenzen: $14.00 - \frac12(16.38 - 2) = 6.81$; en, uit
\ce{Zn(OH)2 + 2OH- <=> Zn(OH)4^2-} ($K = 10^{-1.93}$) met
$[\ce{Zn(OH)4^2-}] = 10^{-2}$, $[\ce{OH-}]^2 = 10^{-0.07}$, pH 13.97.
\ce{Zn^2+}/\ce{Zn}: $-0.76 + 0.030 \times (-2) = -0.82$~V. \ce{Zn(OH)2}/\ce{Zn}:
helling $-0.059$, door $(6.81, -0.82)$: $E = -0.42 - 0.059\,\mathrm{pH}$.
\ce{Zn(OH)4^2- + 4H+ + 2e- -> Zn + 4H2O}: helling $-0.118$, door
$(13.97, -1.24)$: $E = 0.41 - 0.118\,\mathrm{pH}$.
\end{solution}

\begin{solution}{exo:b1:e-ph-diagrams:11}
Het gebied van ijzer ligt bij elke pH onder de waterstoflijn. pH 2:
ijzer wordt geoxideerd tot \ce{Fe^2+} met ontwikkeling van waterstof
(corrosie). pH 7: de reactie is nog mogelijk, maar de waterstoflijn
($-0.41$~V) ligt dicht bij \ce{Fe^2+}/\ce{Fe} en de reactie is zeer
traag. pH 13: ijzer wordt geoxideerd tot \ce{Fe(OH)2}, een
\omterm{def:b1:e-ph-diagrams:corrosion-domains}{passiveringsgebied}: er vormt zich een laag en de spijker blijft blinken.
Met opgeloste zuurstof oxideert de lijn \ce{O2}/\ce{H2O}, die ver erboven
ligt, ijzer bij elke pH tot ijzer(III): roest, \ce{Fe(OH)3} en zijn
gedehydrateerde oxiden.
\end{solution}

\begin{solution}{exo:b1:e-ph-diagrams:12}
Zink heeft het laagste gebied: in contact met het staal wordt het eerst
geoxideerd, \ce{2Zn + O2 + 2H2O -> 2Zn(OH)2} bij pH 8, en de elektronen
die het afgeeft, reduceren de zuurstof aan het staaloppervlak, dat
metallisch blijft. Magnesium ($E^\circ = -2.36$~V) ligt nog lager en werkt
ook; koper (0.34~V) ligt boven ijzer: verbonden met koper zou het ijzer
het metaal zijn dat geoxideerd wordt, sneller dan alleen.
\end{solution}

\begin{solution}{pb:b1:e-ph-diagrams:1}
\textbf{1.} $-$I, 0, $+$I, $+$I.
\textbf{2.} \ce{Cl-} onderaan, \ce{Cl2}, dan \ce{HClO} en \ce{ClO-}
bovenaan.
\textbf{3.} \ce{HClO} en \ce{ClO-}, een zuur-basepaar.
\textbf{4.} \ce{Cl2 + 2e- -> 2Cl-}.
\textbf{5.} \ce{2HClO + 2H+ + 2e- -> Cl2 + 2H2O}.
\textbf{6.} \ce{ClO- + H2O + 2e- -> Cl- + 2OH-}.
\textbf{7.} \ce{HClO} onder pH 7.55, \ce{ClO-} erboven.
\textbf{8.} Er wordt geen elektron uitgewisseld; de grens wordt alleen
door $K_a$ bepaald.
\textbf{9.} $1/(1 + 10^{7.4 - 7.55}) = 0.59$: \qty{59}{\percent}.
\textbf{10.} $1/(1 + 10^{0.45}) = 0.26$: drie kwart van het actieve
chloor zou het zwakkere ontsmettingsmiddel \ce{ClO-} zijn.
\textbf{11.} \ce{ClO-}.
\textbf{12.} $E = 1.396 + 0.0296\log\frac{[\ce{Cl2}]}{[\ce{Cl-}]^2} = 1.396
+ 0.0296\log\frac{0.005}{10^{-4}} = 1.396 + 0.050 = 1.446$~V; er staat
geen \ce{H+} in de \omterm{def:b1:nernst:couple}{halfreactie}.
\textbf{13.} $E = 1.594 + 0.0296\log\frac{[\ce{HClO}]^2[\ce{H+}]^2}{[\ce{Cl2}]}
= 1.594 + 0.0296\log\frac{10^{-4}}{0.005} - 0.0592\,\mathrm{pH} = 1.594 -
0.050 - 0.0592\,\mathrm{pH}$.
\textbf{14.} Gelijkstellen geeft $0.0592\,\mathrm{pH} = 1.594 - 1.396 - 2 \times 0.0503
= 0.0974$, $\mathrm{pH} = 1.646$.
\textbf{15.} \ce{HClO + H+ + 2e- -> Cl- + H2O}: één proton voor twee
elektronen, helling $-0.0592/2 = -0.0296$.
\textbf{16.} Optellen van de twee \omterm{def:b1:nernst:couple}{halfreacties} (elk één elektron per
chlooratoom): $2E^\circ = 1.594 + 1.396$, $E^\circ(\ce{HClO}/\ce{Cl-}) =
1.495$~V; grens $E = 1.495 - 0.0296\,\mathrm{pH}$.
\textbf{17.} \ce{ClO- + 2H+ + 2e- -> Cl- + H2O}: $E = 1.718 - 0.0592\,\mathrm{pH}$
(constante gekozen voor de continuïteit: $1.495 - 0.0296 \times 7.55 = 1.272 = 1.718
- 0.0592 \times 7.55$).
\textbf{18.} \ce{Cl-} onder de lijnen; \ce{Cl2} in een kleine driehoek bij
$\mathrm{pH} < 1.65$ boven 1.446~V; \ce{HClO} boven de lijn met helling
$-0.0296$ tot pH 7.55, \ce{ClO-} daarna. De lijn \ce{O2}/\ce{H2O}, van
1.23 tot 0.40~V, ligt onder al deze grenzen (het berekende diagram van
het script van het hoofdstuk heeft dezelfde vorm).
\textbf{19.} Nee: zijn gebied ligt boven de zuurstoflijn, zodat het water
kan oxideren, maar de reactie is traag. Een zwembad verliest zijn chloor
vooral doordat het oxideert wat zwemmers en de lucht aanbrengen, en door
zonlicht; het moet aangevuld worden.
\textbf{20.} \ce{Cl2} heeft daar geen gebied: het disproportioneert,
\ce{Cl2 + H2O -> HClO + Cl- + H+}; in basisch milieu
\ce{Cl2 + 2OH- -> ClO- + Cl- + H2O}.
\textbf{21.} Door chloor in natronloog te leiden (de reactie van vraag
20).
\textbf{22.} \ce{HClO + Cl- + H+ -> Cl2 + H2O}, een \omterm{def:b1:e-ph-diagrams:disproportionation}{comproportionering}.
\textbf{23.} Onder pH 1.6 comproportioneren hypochloriet en chloride: er
ontstaat chloor, dat boven zijn \omterm{def:b1:precipitation:solubility}{oplosbaarheid} ontsnapt als een giftig
gas.
\textbf{24.} Bij $c = \qty{0.010}{mol/L}$ ligt pH 4 buiten het gebied van
\ce{Cl2}. Maar de grens hangt af van $c$: de vragen 12--14 herhalen voor
een algemene $c$ geeft $\mathrm{pH} = 3.34 + \log(2c)$, wat 3.6 bereikt
voor $c = \qty{1}{mol/L}$, de grootteorde van de concentratie van
bleekwater. Geconcentreerde producten mengen in de buurt van pH 4 is dus
evenmin veilig.
\textbf{25.} \textbf{pH 1.6}: daarboven disproportioneert opgelost chloor
bij $c = \qty{0.010}{mol/L}$ tot hypochlorigzuur en chloride.
\end{solution}
